\documentclass[11pt,a4paper]{article}

\usepackage[utf8]{inputenc}
\usepackage[T1]{fontenc}
\usepackage[french]{babel}
\usepackage{lmodern}

\usepackage{geometry}
\geometry{margin=2.2cm}

\usepackage{hyperref}
\usepackage{booktabs}
\usepackage{longtable}
\usepackage{array}
\usepackage{makecell}
\usepackage{enumitem}
\usepackage{amssymb}

\newcommand{\checkbox}{\(\square\)}

\title{Projet GL --- Documentation de validation\\
Gestion des risques et gestion des rendus}
\author{Groupe 53 : \textit{Les misérables}}
\date{\textit{07-Janvier-2026}}

\begin{document}
\maketitle

\section{Gestion des risques}

\subsection{Objectif}
Un projet logiciel comporte inévitablement des erreurs (interprétation, implémentation, intégration).
L'objectif de cette section est de \textbf{prioriser les dangers critiques} (impact élevé sur la qualité du
projet) et de définir des \textbf{actions concrètes et vérifiables} (réponse oui/non) permettant de réduire
la probabilité d'occurrence et l'impact de ces risques.

\vspace{0.3em}
\noindent \textbf{Rappel :} certains bugs peuvent être rares mais coûteux (compilation cassée).
Nous privilégions donc les risques à \textbf{impact fort}, même si leur correction est ``simple''.

\subsection{Table d'analyse des risques}

% Alignement en haut (p{}) et texte à gauche.
\renewcommand\theadalign{tl}
\renewcommand\cellalign{tl}

\begin{longtable}{|p{1.4cm}|p{4.1cm}|p{1.6cm}|p{1.9cm}|p{6.2cm}|}
\hline
\textbf{ID} & \textbf{Risque / danger} & \textbf{Impact} & \textbf{Proba.} & \textbf{Actions (concrètes, vérifiables)} \\
\hline
\endfirsthead

\hline
\textbf{ID} & \textbf{Risque / danger} & \textbf{Impact} & \textbf{Proba.} & \textbf{Actions (concrètes, vérifiables)} \\
\hline
\endhead

\hline
\endfoot

\hline
\endlastfoot

R1 & Oubli d'une date de suivi ou rendu & Fort & Moyen &
\begin{itemize}[leftmargin=*, nosep]
\item Mettre les 4 dates dans un agenda partagé + rappels
\item Pinner un message ``Deadlines'' dans le canal du groupe.
\end{itemize}
\\ \hline

R2 & Dépôt partagé non compilable (oubli \texttt{git add}, conflit, commit cassé). & Très fort & Moyen &
\begin{itemize}[leftmargin=*, nosep]
\item Interdire le commit direct sur \texttt{main}
\item Exiger \texttt{mvn clean} avant merge.
\item Vérif : on refait un clone dans un dossier vide et ça compile.
\end{itemize}
\\ \hline

R3 & Infrastructure de tests incomplète ou trompeuse (faux positifs). & Fort & Moyen &
\begin{itemize}[leftmargin=*, nosep]
\item Ajouter 3 tests ``sanity'' (valid+invalid) par passe (syntaxe, contexte).
\item Exécuter le script \texttt{common-tests.sh} avant rendu.
\end{itemize}
\\ \hline

R4 & Régression : une correction casse une fonctionnalité déjà validée. & Fort & Moyen &
\begin{itemize}[leftmargin=*, nosep]
\item MR/PR obligatoire + relecture par 1 personne.
\item Vérif : le test reproduisant le bug est ajouté au dépôt.
\end{itemize}
\\ \hline

R5 & Mauvaise interprétation de la grammaire  & Moyen & Moyen &
\begin{itemize}[leftmargin=*, nosep]
\item Identifier règles sensibles (priorités opérateurs, assign, conversions).
\item Créer des tests ciblés (petits) + 1 test ``mixte'' (long).
\end{itemize}
\\ \hline

R6 & Problème environnement (Java/Maven/ANTLR/IMA) & Fort & Moyen &
\begin{itemize}[leftmargin=*, nosep]
\item Vérifier versions : \texttt{java -version}, \texttt{mvn -v}.
\item Vérifier IMA : \texttt{ima -v} sur la machine de rendu.
\item Vérif : exécution d'un exemple minimal end-to-end (compile + test).
\end{itemize}
\\ \hline

R7 & Répartition du travail déséquilibrée / blocage d'une personne. & Moyen & Moyen &
\begin{itemize}[leftmargin=*, nosep]
\item Pair-programming si blocage $>$ 1h.
\item Vérif : chaque tâche a un owner + une date de fin.
\end{itemize}
\\ \hline

R8 & Rendu incomplet : documentation manquante / non à jour. & Fort & Moyen &
\begin{itemize}[leftmargin=*, nosep]
\item Gabarit doc de validation + sections obligatoires
\item Relecture doc la veille du suivi.
\item Vérif : PDF généré depuis le dépôt
\end{itemize}
\\ \hline

\end{longtable}

\subsection{Risques prioritaires}
Nous considérons comme \textbf{prioritaires} les risques à impact très fort (R1, R2, R3, R6) car ils peuvent
rendre le compilateur inutilisable ou décrédibiliser le projet (dépôt cassé, tests non fiables, environnement non opérationnel).

\section{Gestion des rendus (Release management process)}

\subsection{Objectif}
Même avec des tests automatisés, certains contrôles ne sont pas totalement automatisables
(ex : vérifier que le dépôt partagé fonctionne réellement, ou vérifier l'environnement).
Nous définissons donc une \textbf{checklist de rendu} : une suite d'actions simples, vérifiables,
à exécuter avant chaque suivi / rendu.

\subsection{Checklist avant chaque suivi / rendu}
\subsubsection*{A. Vérification dépôt / état Git}
\begin{itemize}[leftmargin=*, itemsep=0.2em]
\item \checkbox\ \texttt{git status} : aucun fichier critique non committé. \textit{(évite l'oubli de fichiers)}
\item \checkbox\ \texttt{git pull} sur la branche cible et résolution propre des conflits. \textit{(évite divergence)}
\end{itemize}

\subsubsection*{B. Build + tests automatiques obligatoires}
\begin{itemize}[leftmargin=*, itemsep=0.2em]
\item \checkbox\ \texttt{mvn -q clean test} \textit{(compile + exécute tests + common-tests.sh)}
\item \checkbox\ Lancer au moins 1 test ``sanity'' par passe :
\begin{itemize}[leftmargin=*, nosep]
\item \checkbox\ syntaxe valid + invalid
\item \checkbox\ contexte valid + invalid
\end{itemize}
\end{itemize}

\subsubsection*{C. Vérification environnement (si codegen/IMA concerné)}
\begin{itemize}[leftmargin=*, itemsep=0.2em]
\item \checkbox\ \texttt{ima -v} fonctionne dans un nouveau terminal. \textit{(vérifie PATH + binaire)}
\item \checkbox\ Compilation + exécution d'un exemple minimal end-to-end (compile + test). \textit{(preuve que tout marche)}
\end{itemize}

\subsubsection*{D. Vérification ``dépôt partagé'' (reproductibilité)}
\begin{itemize}[leftmargin=*, itemsep=0.2em]
\item \checkbox\ Refaire un clone dans un répertoire vide (ou sur une autre machine). \textit{(pas seulement votre clone local)}
\item \checkbox\ Refaire : \texttt{mvn -q clean test}. \textit{(garantit le rendu utilisable par l'enseignant)}
\end{itemize}

\subsubsection*{E. Documentation / livrables}
\begin{itemize}[leftmargin=*, itemsep=0.2em]
\item \checkbox\ README à jour (setup + commandes test). \textit{(rendu exploitable)}
\item \checkbox\ Doc de validation à jour (incluant cette section risques + checklist). \textit{(conforme consignes)}
\end{itemize}

\subsection{Quand exécuter cette checklist ?}
\begin{itemize}[leftmargin=*]
\item \textbf{Avant chaque suivi}
\item \textbf{La veille + le matin du rendu final} (23/01 soutenance)
\item \textbf{Avant toute fusion (MR/PR)} : au minimum parties A et B
\end{itemize}

\end{document}
